\chapter{Context and State of the Art}
\label{ch:context}

\section{Host Institution}

Amen Bank is a privately held Tunisian commercial bank operating under the supervision of the Banque Centrale de Tunisie (BCT). As of 2023, it operates 150 branches across 9 geographic zones, employing 1{,}139 staff of whom 60.5\% are assigned to the branch network. Its client base is composed of individuals (61\% of deposits), private businesses (33\%), and institutional clients (6\%) --- a composition drawn from the audited annual report that directly calibrates the client archetypes of Chapter~\ref{ch:synthetic}.

The internship was hosted by the \textbf{Direction Centrale du Syst\`eme d'Information} (DCSI), within the \emph{D\'epartement des Projets Agence}, which develops and maintains the information systems supporting branch-level operations.

\section{Branch Teller Operations}

The project scope covers the four teller-window (\emph{guichet}) operations that constitute the core transactional activity of a retail branch:

\begin{itemize}
  \item \textbf{Retrait} (cash withdrawal) --- risk exposure includes identity fraud, unusual withdrawal patterns, excessive amounts or frequency.
  \item \textbf{Versement} (cash deposit) --- the primary vector for \emph{structuring} (smurfing): splitting a large sum into multiple deposits below the reporting threshold.
  \item \textbf{Virement} (bank transfer) --- transfers to unknown beneficiaries, duplicates, amounts inconsistent with the client profile.
  \item \textbf{Remise de ch\`eques} (cheque deposit) --- unknown emitters, amounts inconsistent with history, third-party endorsement patterns.
\end{itemize}

Each operation involves \textbf{two actors}: the client initiating it and the employee processing it at the window. Anomalies can originate from either side --- a client structuring deposits, or an employee processing an unusual volume of high-value transactions for a concentrated set of clients. This dual-actor nature motivates the independent client and employee profiling of Chapter~\ref{ch:architecture}.

\section{Regulatory Framework}

Tunisian banks operate under a layered anti-money laundering (LAB) and counter-terrorism financing (FT) framework involving three institutional actors.

\textbf{BCT} issues binding circulars defining internal control obligations, including the obligation to report transactions exceeding 10{,}000~DT to the financial intelligence unit. Circulars 2018-09, 2017-08, and 2021-05 govern LAB/FT internal controls but are restricted documents not publicly available. The system's regulatory rules are therefore calibrated from thresholds extractable from CTAF case studies and published legislation rather than a complete circular inventory --- a gap documented explicitly rather than papered over.

\textbf{CTAF} (Commission Tunisienne des Analyses Financi\`eres) is Tunisia's financial intelligence unit. Its 2021 annual report provides the behavioral ground truth for what suspicious operations look like locally: \emph{esp\`eces} are involved in 68\% of flagged cases, \emph{virements} in 56\%, \emph{ch\`eques} in 34\%. The concentration of reported amounts around the 10{,}000~DT threshold --- 85\% of flagged operations in the 5{,}000--20{,}000~DT band --- directly informed the near-threshold detection logic of the Decision Service.

\textbf{Law n\textsuperscript{o}41-2024}, effective February 2025, caps individual cheques at 30{,}000~DT. BCT data from Q1 2026 confirms a 24.9\% collapse in cheque usage volume following implementation. This required adjusting the cheque archetype parameters and introduced a new anomaly scenario: cheque amounts clustering just below the ceiling, analogous to structuring below the reporting threshold.

\section{Existing Detection Capabilities and Project Positioning}

The Amen Bank 2023 annual report confirms that the bank already operates an AI-based AML system. Based on the report's description, that system performs \textbf{bank-wide compliance reporting}: a compliance-officer audience, a quarterly tempo, and output aligned with CTAF regulatory requirements.

The system built in this internship occupies a different position under the same regulatory umbrella: \textbf{real-time decision support for branch supervisors} on four specific \emph{caisse} operations, profiling client and employee as independent actors, at near-real-time latency, with sequence-aware LSTM modeling of behavioral drift. The two systems are complementary rather than competing --- the existing one answers \emph{``should the bank file a regulatory report?''} while this project answers \emph{``should the branch supervisor investigate this transaction now?''}

\section{State of the Art}

\subsection{Rule-Based and Statistical Detection}

The earliest and still most widely deployed approach to transaction monitoring relies on deterministic rules: fixed thresholds, velocity checks, and pattern matching defined by compliance officers. Bolton and Hand~\cite{bolton} note that rule-based systems dominate production deployments due to their interpretability and auditability, but have two structural limitations: they require manual maintenance for every new fraud typology, and they cannot detect patterns outside the predefined rule set. Conversely, regulatory requirements often \emph{mandate} deterministic checks (such as the 10{,}000~DT reporting threshold) that cannot be delegated to a probabilistic model. The system therefore retains deterministic rules in the Decision Service while delegating pattern discovery to the deep learning models.

Statistical methods model the distribution of normal behavior and flag significant deviations. \textbf{Isolation Forest}~\cite{iforest} builds an ensemble of random trees in which anomalies, being few and different, require fewer partitions to isolate. \textbf{Local Outlier Factor}~\cite{lof} compares the local density of an observation to that of its neighbors. Both are efficient and distribution-free, and both are fundamentally \emph{point-wise}: they evaluate observations independently and cannot model temporal dependencies. Detecting behavioral drift --- a client whose pattern gradually shifts over weeks --- requires models that consume sequences.

\subsection{Autoencoders}

An Autoencoder is a neural network trained to reconstruct its input through a compressed latent representation. Trained exclusively on normal data, it learns to encode and decode the patterns of legitimate transactions; anomalous inputs produce high reconstruction errors that serve as anomaly scores. Sakurada and Yairi~\cite{sakurada} showed that Autoencoders outperform linear PCA-based methods when normal data occupies a nonlinear manifold, and An and Cho~\cite{ancho} formalized reconstruction error as an anomaly score.

In the financial domain this addresses the class imbalance problem directly: anomalies represent less than 1\% of transactions, making supervised classification impractical without extensive labels. Training only on normal transactions removes the need for labeled anomalies entirely --- a critical advantage where confirmed fraud labels are scarce and delayed by months. The system uses an Autoencoder as its \textbf{point anomaly detector}.

\subsection{LSTM Networks}

LSTM networks~\cite{lstm} address the vanishing gradient problem that prevents standard recurrent networks from learning long-range dependencies; their gating mechanism allows selective retention of information across long sequences. For anomaly detection the LSTM is trained as a \emph{next-step predictor}: given a sequence of recent transactions it predicts the features of the next one, and the deviation between predicted and observed features becomes the anomaly score. Bontemps et al.~\cite{bontemps} applied this to network intrusion detection, showing that it detects subtle sequential anomalies that point-wise methods miss.

In banking, the LSTM captures patterns that unfold over time: a client who normally transacts once a month suddenly executing five transactions in three days, or a gradual increase in transfer amounts to a new beneficiary. These are invisible to the Autoencoder, which evaluates each transaction in isolation. The system uses the LSTM as its \textbf{sequential anomaly detector}.

\subsection{Score Fusion}

When multiple detectors operate in parallel their scores must be combined. Aggarwal~\cite{aggarwal} surveys fusion strategies from simple averaging to learned combination; the key challenge is that different detectors have different score distributions, making naive averaging misleading. The system uses a \emph{sigmoid-weighted} fusion that adapts to account maturity: new accounts favor the Autoencoder (which needs only one event) over the LSTM (which needs a populated sequence buffer), and the LSTM weight increases as the account matures. This dissolves the cold-start problem inherent to sequential models without a separate detection pathway.

\subsection{Explainability with SHAP}

Deploying a deep model in a regulated environment requires that every alert carry an explanation a non-technical supervisor can act on. SHAP~\cite{shap} attributes a prediction to individual input features using game-theoretic Shapley values, with three properties that suit this context: local accuracy, missingness, and consistency.

For the Autoencoder the system uses \textbf{KernelSHAP}, a model-agnostic approximation treating the reconstruction-error function as a black box. DeepSHAP is inapplicable because the anomaly score subtracts the input from the reconstruction, breaking the forward-pass tracing DeepSHAP requires. For the LSTM, SHAP values over the full sequence-length-by-feature tensor are too granular for a supervisor to read; the system instead reports an \emph{expected-versus-actual} comparison per interpretable feature, making deviations immediately legible.

\subsection{Streaming Architectures}

Real-time detection requires processing events as they arrive rather than in periodic batches. Kreps et al.~\cite{kafka} introduced the log-based streaming paradigm, where events are written to an append-only log partitioned by key and consumed by independent services at their own pace, decoupling producers from consumers and enabling horizontal scaling of individual stages. The system adopts this paradigm using RedPanda, a Kafka-compatible broker that eliminates the JVM and ZooKeeper dependencies while preserving API compatibility. Partitioning by client identifier ensures that all transactions for a given client are processed in order --- a requirement of the LSTM's sequential modeling.

\section{Problem Statement and Objectives}

Detection of anomalous teller operations at Amen Bank currently relies on manual rules and periodic compliance reviews. This cannot scale to the transaction volume of 150 branches, cannot adapt to evolving typologies without manual rule updates, and provides no real-time decision support to branch supervisors.

The objective of this internship is to design and implement a real-time anomaly detection platform that:

\begin{enumerate}
  \item ingests teller transactions as they occur and processes them through a streaming pipeline at near-real-time latency;
  \item profiles both the client and the employee as independent behavioral actors;
  \item combines a point anomaly detector (Autoencoder) with a sequential one (LSTM) to capture both individual deviations and behavioral drift;
  \item accompanies every alert with a human-readable SHAP explanation;
  \item provides a supervisor dashboard for alert review, operational monitoring, and transaction simulation;
  \item implements the MLOps infrastructure --- experiment tracking, automated retraining, drift detection, observability --- required to maintain the system beyond initial deployment.
\end{enumerate}
